% Осень26!A27; нумерация внутри столбца без пустых строк.
% Источник: лекция 4, §4 (строка 106)
% Источник: лекция 4, §5 (строка 158)
% Источник: лекция 4, §7 (строка 226)
\ticket{24}{Регулярность меры Лебега}
\label{ticket:24}

\subsection*{Теорема о регулярности меры Лебега}

Свойство «хорошей приспособленности» меры Лебега к структуре открытых/замкнутых множеств \(\mathbb{R}^m\).

\begin{theorem}

\emph{регулярность меры Лебега.} Пусть \(A\in\mathcal{M}^m\) (измеримо по Лебегу). Тогда для любого \(\varepsilon>0\):

\begin{enumerate}
\def\labelenumi{\arabic{enumi}.}
\tightlist
\item
  существует \textbf{открытое} множество \(G_\varepsilon\supset A\) такое, что \(\lambda(G_\varepsilon\setminus A)<\varepsilon\) (аппроксимация снаружи);
\item
  существует \textbf{замкнутое} множество \(F_\varepsilon\subset A\) такое, что \(\lambda(A\setminus F_\varepsilon)<\varepsilon\) (аппроксимация изнутри).
\end{enumerate}

Смысл: «изучая любое измеримое множество \(A\), с точки зрения теории меры можно, незначительно поменяв меру, подменить его на чуть более широкое открытое или на чуть меньшее замкнутое».

\end{theorem}

\subsubsection*{\texorpdfstring{Доказательство пункта 1: случай \(\lambda A<+\infty\)}{Доказательство пункта 1: случай \textbackslash lambda A\textless+\textbackslash infty}}

По п. 5 теоремы о лебеговском продолжении меры \[\lambda A=\inf\Bigl\{\sum_k\lambda P_k:\ A\subset\bigcup_k P_k,\ P_k\ \text{— ячейки}\Bigr\}.\]

\emph{Замечание о конечности (важно).} Если \(\lambda A<+\infty\), то покрытие ячейками можно выбрать «плотно облегающим» \(A\): по полуаддитивности \(\lambda A\le\lambda\bigl(\bigcup_kP_k\bigr)\le\sum_k\lambda P_k\), и, оптимизируя ячейки, зазор \(\lambda\bigl(\bigcup_kP_k\bigr)-\lambda A\) делается сколь угодно малым. Если же \(\lambda A=+\infty\), формула инфимума «не несёт информации» о плотности облегания. Это же объясняет, почему в теореме о продолжении нужна была \textbf{σ-конечность}: на множествах конечной меры возможны такие аппроксимации, а на бесконечной «всё теряется». Мы в \(\mathbb{R}^m\) сначала рассмотрим конечные куски, потом воспользуемся σ-конечностью.

Итак, для данного \(\varepsilon>0\) найдутся ячейки \(P_k\) такие, что \[A\subset\bigcup_kP_k,\qquad \lambda\Bigl(\bigcup_k P_k\setminus A\Bigr)<\frac{\varepsilon}{2}.\]

Теперь «разденем/раздуем» ячейки: каждая ячейка \(P_k=[a_k,b_k)\) --- «ни то ни сё», полуоткрытая; превратим её в \textbf{открытый} параллелепипед \(A_kB_k=(a_k',b_k)\supset P_k\), чуть уменьшив все координаты вектора \(a_k\), с контролем меры зазора: \[\lambda\bigl((a_k',b_k)\setminus P_k\bigr)<\frac{\varepsilon}{2\cdot 2^k}\] (запас \(\varepsilon/2\) на итог и \(2^{-k}\), чтобы ряд свернулся при суммировании).

Положим \(G_\varepsilon:=\bigcup_k(a_k',b_k)\) --- открытое множество, \(G_\varepsilon\supset A\). Тогда

\[G_\varepsilon\setminus A\;\subset\;\Bigl(\bigcup_kP_k\setminus A\Bigr)\;\cup\;\bigcup_k\bigl((a_k',b_k)\setminus P_k\bigr),\]

--- «наивная теория множеств, ясная кому угодно» (из объединения «матрёшек» \((a_k',b_k)\supset P_k\) выбрасывается \(A\); всё, что не вылетело в зазор «кольцо между \(A_kB_k\) и \(P_k\)», осталось внутри \(\bigcup P_k\)). Отсюда по полуаддитивности и аддитивности (для вложенных измеримых множеств):

\[\lambda(G_\varepsilon\setminus A)\;\le\;\lambda\Bigl(\bigcup_kP_k\setminus A\Bigr)+\sum_k\lambda\bigl((a_k',b_k)\setminus P_k\bigr)<\frac\varepsilon2+\sum_k\frac{\varepsilon}{2\cdot2^k}=\frac\varepsilon2+\frac\varepsilon2=\varepsilon.\qquad\blacksquare\]

\subsubsection*{\texorpdfstring{Доказательство пункта 1: случай \(\lambda A=+\infty\) (общий случай)}{Доказательство пункта 1: случай \textbackslash lambda A=+\textbackslash infty (общий случай)}}

Смотрим «в тетрадку»: \(\mathbb{R}^m\) разбито в дизъюнктное объединение ячеек («клеточки»), \[\mathbb{R}^m=\bigsqcup_{k=1}^\infty Q_k,\qquad \lambda Q_k<+\infty\] (например, целочисленные кубы). Рассмотрим \(A_k:=A\cap Q_k\) --- измеримо (пересечение измеримых), \(\lambda A_k<+\infty\). К каждому \(A_k\) применим случай 4.1 с точностью \(\varepsilon/2^k\): найдётся открытое \(G^{(k)}\supset A_k\) с \(\lambda(G^{(k)}\setminus A_k)<\varepsilon/2^k\). Тогда \[G_\varepsilon:=\bigcup_kG^{(k)}\] открыто, содержит \(A\), и \[\lambda(G_\varepsilon\setminus A)\;\le\;\sum_k\lambda\bigl(G^{(k)}\setminus A\bigr)\;\le\;\sum_k\lambda\bigl(G^{(k)}\setminus A_k\bigr)<\sum_k\frac{\varepsilon}{2^k}=\varepsilon\] («при подсчёте зазоров с нахлёстом получится ещё меньше \(\varepsilon\)» --- лектор признал, что без формул это выглядит как «махание руками», но «совершенно ясно»). \(\blacksquare\)

\subsubsection*{Доказательство пункта 2 (замкнутые множества) --- переход к дополнению}

«Чисто организационная загадка, думать не надо». Рассмотрим дополнение \(A^c:=\mathbb{R}^m\setminus A\in\mathcal{M}^m\). По пункту 1 найдётся открытое \(G\supset A^c\) с \(\lambda(G\setminus A^c)<\varepsilon\). Положим \[F_\varepsilon:=G^c\;(=\mathbb{R}^m\setminus G).\] \(F_\varepsilon\) замкнуто (дополнение открытого) и \(F_\varepsilon\subset A\). Картинка: «овальная рамочка» между пунктирным овалом \(G\) (внешним для \(A^c\)) и самим \(A^c\) --- та же самая, что и зазор между \(F_\varepsilon\) и \(A\): \[A\setminus F_\varepsilon=A\cap G=G\setminus A^c\quad\Longrightarrow\quad \lambda(A\setminus F_\varepsilon)=\lambda(G\setminus A^c)<\varepsilon.\qquad\blacksquare\]

\subsection*{Следствия регулярности: борелевская аппроксимация; множества меры ноль; примеры}

\subsubsection*{Следствие 1: зажатие между борелевскими множествами}

Здесь «сыграло» определение борелевской σ-алгебры из §3.

\textbf{Следствие 1.} Пусть \(A\in\mathcal{M}^m\). Тогда существуют борелевские множества \(B\subset A\subset C\) (т.е. \(B,C\in\mathcal{B}(\mathbb{R}^m)\)) такие, что \[\lambda(C\setminus B)=0.\] В частности, \(\lambda(A\setminus B)=0\) и \(\lambda(C\setminus A)=0\) (оба множества --- подмножества \(C\setminus B\)).

\begin{proof}

В качестве \(C\) возьмём счётное пересечение открытых приближений из п. 1 теоремы: \[C:=\bigcap_{n=1}^\infty G_{1/n},\] где \(G_{1/n}\supset A\) открыто и \(\lambda(G_{1/n}\setminus A)<1/n\). Тогда \(C\) борелевское, \(C\supset A\), и при каждом \(n\) \[C\setminus A\subset G_{1/n}\setminus A\quad\Rightarrow\quad\lambda(C\setminus A)<\frac1n,\] значит \(\lambda(C\setminus A)=0\). Аналогично \(B:=\bigcup_{n=1}^\infty F_{1/n}\) (счётное объединение замкнутых) борелевское, \(B\subset A\), \(\lambda(A\setminus B)=0\). Тогда \(C\setminus B=(C\setminus A)\cup(A\setminus B)\) имеет меру ноль.

\end{proof}

\textbf{Интерпретация.} Если «пренебрегать множествами меры ноль», всякое измеримое множество можно считать борелевским: к \(A\) можно чуть-чуть добавить множества меры ноль (получится борелевское \(C\supset A\)) или чуть-чуть выбросить (получится борелевское \(B\subset A\)) --- «по мере ничего не потеряется, а множество станет простым по структуре».

\textbf{Контринтуитивный пример (рациональные точки квадрата).} Пусть \(A=(\mathbb{Q}\cap[0,1])^2\subset\mathbb{R}^2\) --- рациональные точки единичного квадрата. \(A\) счётно, \(\lambda A=0\). По п. 1 теоремы о регулярности все рациональные точки квадрата накрываются \textbf{открытым} множеством \(G\) меры \(<1/1000\) (или любого наперёд заданного \(\varepsilon>0\)). «Это точно соответствует вашей геометрической интуиции? Рациональные точки расположены густо, а накрыть их надо открытыми параллелепипедами суммарной меры \(1/1000\)». Это доказано, но требует осознания и принятия.

\subsubsection*{Следствие 2: структура измеримого множества}

\textbf{Следствие 2.} Всякое \(A\in\mathcal{M}^m\) представимо в виде \[A=B\cup N,\] где \(B\) --- борелевское множество, \(N\) --- множество меры ноль (\(\lambda N=0\)).

\begin{proof}

\(B\) --- из следствия 1; \(N:=A\setminus B\); по следствию 1 \(\lambda N=0\).

\end{proof}

Каждое измеримое множество --- это «борелевское (то, что строится из открытых) плюс добавка меры ноль», а мера ноль --- «это можно накрыть сколь угодно малой системой ячеек».

\subsubsection*{Критерий меры ноль (без априорной измеримости)}

Лектор дополнил следствия важным наблюдением в обе стороны.

\textbf{Предложение.} Пусть \(N\subset\mathbb{R}^m\) обладает свойством: \[\forall\varepsilon>0\ \exists\ \text{ячейки } P_k^{(\varepsilon)}:\quad N\subset\bigcup_{k=1}^\infty P_k^{(\varepsilon)},\qquad \sum_k\lambda P_k^{(\varepsilon)}<\varepsilon.\] Тогда \(N\) измеримо по Лебегу и \(\lambda N=0\).

\begin{proof}

Для каждого \(l\in\mathbb{N}\) возьмём покрытие ячейками \(N\subset\bigcup_kP_k^{(l)}\) с \(\sum_k\lambda P_k^{(l)}<1/l\). Рассмотрим \[M:=\bigcap_{l=1}^\infty\Bigl(\bigcup_kP_k^{(l)}\Bigr).\] \(M\) измеримо и даже борелевское (ячейка --- борелевское множество: она получается из открытых параллелепипедов предельным переходом, «раздуя-сдувая»); \(N\subset M\). По монотонности и полуаддитивности \[\lambda M\le\lambda\Bigl(\bigcup_kP_k^{(l)}\Bigr)\le\sum_k\lambda P_k^{(l)}<\frac1l\qquad\forall l,\] следовательно \(\lambda M=0\). Итак, \(N\) --- подмножество множества \(M\) меры ноль. Вступает «тяжёлая артиллерия» --- \textbf{полнота} меры Лебега: всякое подмножество множества меры ноль измеримо и имеет меру ноль. Значит, \(N\in\mathcal{M}^m\), \(\lambda N=0\).

\end{proof}

\textbf{Итог:} «множество измеримо и имеет меру ноль \(\iff\) его можно накрывать сколь угодно малыми (по суммарной мере) системами ячеек» --- свойство меры ноль легко проверяется.

\subsubsection*{Канторово множество под новым углом}

\textbf{Пример (канторовское множество).} Из предыдущей лекции: \[K=\Bigl\{x\in[0,1]:\ x=0{,}\alpha_1\alpha_2\alpha_3\ldots \text{ в троичной системе, все } \alpha_i\in\{0,2\}\Bigr\}.\]

\begin{itemize}
\tightlist
\item
  Таких чисел \textbf{континуум}: каждая последовательность нулей и двоек даёт точку \(K\) (взаимно однозначно, с оговоркой о двояких представлениях).
\item
  Геометрия: первая троичная цифра не 1 --- выбрасываем среднюю треть, и т.д. (узнаваемая конструкция Кантора).
\item
  На \(n\)-м шаге \(K\) содержится в объединении \(2^n\) отрезков суммарной длины \((2/3)^n\): было 1, потом \(2/3\), потом \((2/3)^2\), \ldots{}
\item
  \(\forall\varepsilon>0\) при больших \(n\): \((2/3)^n<\varepsilon\). По предложению 5.3 \(K\) \textbf{измеримо и \(\lambda K=0\)} --- даже если «не знать» заранее об измеримости.
\end{itemize}

\subsection*{Следствие 3: мера через открытые, замкнутые и компактные множества}

«Сводное утверждение, пересказ проделанной работы на третий язык» --- оправдание того, что сложный рецепт теоремы о продолжении (инфимум по всем счётным покрытиям ячейками) можно заменить простым «геометрическим» перебором.

\textbf{Следствие 3.} Для всякого \(A\in\mathcal{M}^m\): \[
\begin{aligned}
\lambda A
 &=\inf\{\lambda G:\ G\ \text{открыто},\ A\subset G\}\\
 &=\sup\{\lambda F:\ F\ \text{замкнуто},\ F\subset A\}\\
 &=\sup\{\lambda K:\ K\ \text{компактно},\ K\subset A\}.
\end{aligned}
\]

Первые два равенства --- прямое следствие теоремы о регулярности (п. 1 даёт \(\lambda A\le\lambda G_\varepsilon<\lambda A+\varepsilon\); п. 2 --- аналогично снизу). Новое --- третье (компакты; в \(\mathbb{R}^m\) компакт = замкнуто + ограниченно).

\textbf{Схема доказательства третьего равенства.}

\emph{Случай ограниченного \(A\):} всякое замкнутое подмножество ограниченного множества ограничено, т.е. компактно; перебирая замкнутые подмножества \(A\), мы «реально перебираем компакты» --- равенство «звёздочка очевидна».

\emph{Случай неограниченного \(A\):} разобьём \(\mathbb{R}^m=\bigsqcup_{k=1}^\infty Q_k\) на ячейки конечной меры («клеточки», можно целочисленные) и положим \[A_n:=\bigcup_{k=1}^{n}(A\cap Q_k).\] Тогда \(A_1\subset A_2\subset\cdots\), \(A=\bigcup_nA_n\), и по \textbf{непрерывности меры снизу} \[\lambda A=\lim_{n\to\infty}\lambda A_n=\lim_{n\to\infty}\sum_{k=1}^{n}\lambda(A\cap Q_k).\]

\begin{itemize}
\tightlist
\item
  \textbf{Если \(\lambda A<+\infty\):} подберём \(N\) так, что \(\lambda A-\lambda A_N<\varepsilon/2\) («откинув хвост объединения, по мере потеряем меньше \(\varepsilon/2\)»). Множество \(A_N\) ограничено (конечное объединение ограниченных ячеек) и конечной меры; по п. 2 теоремы найдётся замкнутое \(K\subset A_N\) с \(\lambda(A_N\setminus K)<\varepsilon/2\); будучи замкнутым подмножеством ограниченного \(A_N\), \(K\) \textbf{компактно}. Тогда \[\lambda(A\setminus K)\le\lambda(A\setminus A_N)+\lambda(A_N\setminus K)<\frac\varepsilon2+\frac\varepsilon2=\varepsilon.\]
\item
  \textbf{Если \(\lambda A=+\infty\):} для любого большого \(E>0\) подберём \(N\) с \(\lambda A_N>E\); внутри \(A_N\) (ограниченного, конечной меры) подберём компакт \(K\subset A_N\), «лишь чуть-чуть проигрывающий по мере» этому объединению --- так, что в зазор между \(\lambda A_N\) и \(E\) «втискивается» мера компакта, т.е. \(\lambda K>E\). Отсюда \(\sup\{\lambda K:\ K\subset A\ \text{компактно}\}=+\infty=\lambda A\). \(\blacksquare\)
\end{itemize}

\begin{definition}

\emph{регулярность.} Совокупность свойств теоремы §4 и следствий 1--3 --- выраженность меры множества через меры открытых, замкнутых и компактных множеств --- называется \textbf{регулярностью меры Лебега}.

\end{definition}
